\section{The framework of Alman and Vassilevska Williams}
\label{sec:framework}

We recall the parts of \cite{AV} that are used.  Numbers of lemmas and theorems
``of the paper'' are those of \cite{AV}.  All logarithms are natural unless a
base is written, \(\Hent(x)=-x\ln x-(1-x)\ln(1-x)\) is the entropy function, and
\(\psi(x)=x\ln x\) with \(\psi(0)=0\).

\subsection{Problems and the machine}
\label{sec:machine}

\emph{Exact Triangle} asks whether a weighted tripartite graph with \(n\)
vertices per part has a triangle of weight zero.  \emph{3SUM} asks whether \(n\)
integers contain three with sum zero.  The \emph{\(\minplus\)-product} of two
\(n\times n\) integer matrices \(A,B\) is \(C_{ij}=\min_k(A_{ik}+B_{kj})\), and
\emph{APSP} asks for all distances in a directed graph with integer weights and
no negative cycle.

The \emph{thin product} is the following task.  Given \(X\in\Z^{N\times D}\),
\(Y\in\Z^{D\times N}\) and a set \(W\) of \emph{wanted} positions of the
\(N\times N\) matrix \(XY\), compute \((XY)[I,J]\) for all \((I,J)\in W\).  The
number \(D\) is the \emph{inner dimension}.  \(\LopTri(n,D)\)
(Definition~13 of the paper) is the special case that Exact Triangle reduces
to: a tripartite graph with outer parts of \(n\) vertices and a middle part of
at most \(D\) vertices, and a set \(W\) of \emph{query pairs} of outer
vertices; decide for each pair whether it has a common neighbor in the middle
part.

Running times are those of programs of a word RAM, in the sense of the
formalization \cite{upstream} of the paper.  For a problem \(Q\) and a rational
number \(r\), the statement \lean{Q.SolvedInTime}~\(r\) says: for every \(\kappa\)
there are a program, a slope \(b\) and a bound \(T(n)=O(n^{r})\) such that the
program is correct, within \(T(n)\) steps, on all instances of size \(n\) whose
numbers have absolute value at most \(n^{\kappa}\), at every word size
\(\ge b(\lfloor\log_2 n\rfloor+1)\).  The formalization proves Theorem~22 of
the paper in this form, with \(r=\ThreeSumPaperStated\) for 3SUM and
\(r=\APSPpaper\) for the \(\minplus\)-product and for APSP.

\subsection{The recursion}
\label{sec:recursion}

The recursion of Section~2 of the paper applies Sch\"onhage's identity
\cite{Sch81} for \(\schon313\oplus\schon141\).  The identity has ten terms.
Nine of them feed one outer output each, and the tenth, \(\Pz\), is shared.  It
has seven left variables: three outer ones and four inner ones.

The identity is applied on \(L\) levels, \(m\) of which are inner.  We write
\[
  D=4^{m},\qquad N_0=3^{L-m},\qquad K=\binomial Lm,\qquad
  M=K N_0^{2}=\binomial Lm 9^{L-m},\qquad c=\frac Lm\ge 10 .
\]
The matrix \(XY\) is cut into \emph{tiles} of side \(\sqrt K N_0\) (at most
\(4N^2/M\) of them), each with \(M\) output strings, and the rows of \(X\) into
\emph{bands} of \(\sqrt K N_0\) rows.

A \emph{leaf} is a string of \(L\) terms.  Its \emph{order} is \(m\) minus its
number of symbols \(\Pz\).  For a set \(U\) of output strings, \(\Leaves(U)\)
is the set of leaves at which the recursion multiplies for some string of
\(U\); these leaves have order between \(0\) and \(m\).  Two counts govern the
cost:
\[
  \alpha_d=\binomial md 9^{d},\qquad
  \beta_d=\binomial L{m-d}9^{L-m+d}\qquad(0\le d\le m).
\]
One output string is fed by \(\alpha_d\) leaves of order \(d\), and a tile has
\(\beta_d\) leaves of order \(d\); \(\beta_0=M\).  Lemma~11 of the paper is the
bound
\begin{equation}\label{eq:lemma11}
  |\Leaves(U)|\le\sum_{d=0}^{m}\min\bigl(|U|\alpha_d,\ \beta_d\bigr)
\end{equation}
for a set \(U\) of output strings of one tile, followed by an estimate of the
right-hand side.

\subsection{Boxes and queries}
\label{sec:boxes}

Section~4 of the paper (Theorem~30) turns the recursion into a data structure
with a switching order \(0\le t\le m\).  The leaves of order at least \(t\) are
covered by precomputed \emph{boxes}, at most
\begin{equation}\label{eq:boxes}
  (m+1)\sum_{d=t}^{m}\beta_d
\end{equation}
per tile (Lemma~29 of the paper).  A query for one entry sums over the leaves
of order below \(t\) and costs \(O\bigl(L\sum_{d\le t}\alpha_d\bigr)\).  The
encodings \(\Phi_\tau(a)\) of the input at all \(10^{L}\) leaves \(\tau\) are
computed for every band by Yates' algorithm \cite{Yat37}.  The preprocessing
costs a constant times the expression~(8) of the paper,
\begin{equation}\label{eq:cost8}
  L\,m\,\frac{\rho^{t}}{1-\rho}\,N^{2}+\frac{N\cdot 10^{L}}{\sqrt K N_0},
  \qquad \rho=\frac{9m}{L-m+1},
\end{equation}
because the paper bounds \(\sum_{d\ge t}\beta_d\le M\rho^{t}/(1-\rho)\); the
second summand is the cost of the encodings.

Fix real numbers \(c>10\) and \(0<\theta<0.9\) and let \(L=\lceil cm\rceil\),
\(t=\lceil\theta m\rceil\).  Then \(\rho<\rho_c:=9/(c-1)<1\).  Put
\begin{equation}\label{eq:gpap-q}
  \gpap(c,\theta)=\frac{\theta\ln(1/\rho_c)}{\ln 4},\qquad
  q(\theta)=\frac{\Hent(\theta)+\theta\ln 9}{\ln 4}.
\end{equation}
The first summand of \eqref{eq:cost8} is \(O(N^{2}\log^{2}D/D^{\gpap})\) and a
query costs \(O(D^{q(\theta)}\log D)\).  The second summand is at most
\(N^{2}/D^{\gamma}\) when the instance is \emph{thin}: \(D\le N^{\varepsilon}\)
with
\begin{equation}\label{eq:Rc}
  \varepsilon<\Rc(\gamma):=\frac{\ln 4}{\Afull(c)+\gamma\ln 4},\qquad
  \Afull(c)=c\ln 10-\frac c2\Hent\Bigl(\frac1c\Bigr)-(c-1)\ln 3 .
\end{equation}
This is equation~(11) of the paper; \(\Afull(c)\) is the exponential rate of
\(10^{L}/(\sqrt K N_0)\) per inner level.  Corollary~31 of the paper is the
resulting data structure, and Corollary~32 the offline form: for
\(|W|\le N^{2}/D^{\kappa}\) wanted entries the thin product takes time
\begin{equation}\label{eq:cor32}
  O\bigl(N^{2}\log^{2}D\,(D^{-\gamma}+D^{q(\theta)-\kappa})\bigr),
  \qquad\gamma=\gpap(c,\theta).
\end{equation}
Corollary~26 of the paper takes \(c=21\), \(\theta=1/9\), \(\kappa=1/2\) and
\(\varepsilon=1/18\), which gives \(\gamma=0.063\).

\subsection{From the thin product to Exact Triangle}
\label{sec:host}

Theorem~17 of the paper reduces Exact Triangle on \(n\) vertices per part with
weights of absolute value at most \(n^{\nu}\) to \(\LopTri\).  Let
\(16\le D'\le n\) and let \(1\le g'\le\sqrt{D'}\) be an integer.  The weights
are hashed modulo a prime \(p\in[\sqrt{D'}/2,\sqrt{D'})\) that is chosen to
minimize the number of false positives.  The reduction produces at most
\(4ng'\) instances of \(\LopTri(n,D')\) with at most \(n^{2}/\sqrt{D'}\) query
pairs each, and takes
\begin{equation}\label{eq:thm17}
  O\Bigl(\frac{\nu n^{3}\log n}{g'}+n^{\omega+o(1)}D'^{3/2}+n^{2}D'g'\Bigr)
\end{equation}
additional time.  The three terms are the scans for witnesses, the choice of the
prime and the construction of the instances.  The middle part of an instance is
\(C_k\times\Z_p\) for a piece \(C_k\) of one part with
\(|C_k|\le\lceil s/g'\rceil\), \(s=\lfloor\sqrt{D'}\rfloor\).  We write \(D'\)
and \(g'\) for the parameters of the reduction and keep \(D\) for the inner
dimension of the thin products; in the paper the two agree.

With \(D'=n^{1/18}\) and \(g'=D'^{\gamma/2}\), \(\gamma=0.063\), the calls and
the scans balance (Remark~20 of the paper), and Exact Triangle takes time
\(n^{3-\ETpaper}\) up to a logarithmic factor (Theorem~19 of the paper), since
\(0.063/36=\ETpaper\).

\subsection{From Exact Triangle to 3SUM and APSP}
\label{sec:known}

Theorem~21 of the paper recalls two known deterministic reductions.  3SUM on
\(n\) numbers reduces to \(n^{1/2+o(1)}\) instances of Exact Triangle on
\(n^{1/2+o(1)}\) vertices per part \cite{CH20,VW13}.  The \(\minplus\)-product
of \(n\times n\) matrices with entries up to \(U\) takes time
\(O(n^{2}T(n^{1/3})\log^{2}U)\) if Exact Triangle on \(s\) vertices per part
takes time \(T(s)\), and APSP follows with one more logarithm
\cite{VW10,VW18,VW13}.  So if Exact Triangle is solved in time \(n^{3-\delta}\),
then 3SUM is solved in time \(n^{a}\) for every \(a>2-\delta/2\), and the
\(\minplus\)-product and APSP in time \(n^{a}\) for every \(a>3-\delta/3\):
3SUM keeps half of the saving and APSP a third.  With \(\delta=\ETpaper\) this
is Theorem~22 of the paper: 3SUM in time \(n^{\ThreeSumPaper+o(1)}\), printed
as \(O(n^{\ThreeSumPaperStated})\), and APSP in time \(O(n^{\APSPpaper})\).
